\section{敏感性分析：能力收缩与方案阈值}
\subsection{储备变化为何触发组批跳变}
储备比例直接改变可用任务能量，而不改变机型额定载荷。为分离连续与离散效应，对 10\%--40\% 储备设置 61 个等距点，逐一重算 15 服务区、3 机型的载荷边界；另在 10\%、15\%、20\%、25\%、30\% 五个场景重新运行精确组批。
\paperfig{fig07_sensitivity}{储备增加先收缩最受限服务区的载荷边界，再引发组批架次数变化。阴影表示不同服务区的能力范围，并非统计置信区间；不可达状态显示为零交付能力，不与空载可达混同。下图五个点均重新优化，虚线为基准固定组批的最大共同储备 \QOneReserveLimit\%。}{sensitivity}
\papertable{sensitivity}{返航储备变化下的 Q1 精确组批结果}{sensitivity}

基准组批在 10\%--20\% 的三个已计算场景保持相同目标值；储备提高至 25\% 和 30\% 时，架次数分别变为 19 和 20。其原因是单批货箱不可连续缩小：一旦某一批的固定质量越过安全边界，就必须重分配或增加架次。因此连续的能量预算变化可以引起离散的方案跳变。

基准固定方案最小返航 SOC 为 \QOneReserveLimit\%，故在不改变任何货箱与机型时，它只能支持不超过该值的统一储备。该阈值不等于“所有 18 架次方案都失效”的阈值；后者须对临界储备重新求解才能确定，本文不以五个场景间的连线代替未计算的最优解。

\subsection{固定联合计划的能量与通信压力边界}
对于统一能量上浮系数 $\lambda$，在路径、时间和载荷不变时，所有运输架次满足储备的最大系数为
\begin{equation}
\lambda_{\max}=\min_{p\in\mathcal P}
\frac{(1-\rho_{g_p})E_{g_p}}{E_p}.
\label{eq:stress}
\end{equation}
Q2 与 Q3 固定运输计划的 \(\lambda_{\max}-1\) 仅为 \QThreeEnergyTolerance\%。这是一项可解析复核的固定计划压力检查，既不是实测误差分布，也不是已重新优化后的恢复能力。它说明只根据全部架次可行，不能推断方案对实际能耗误差具有较强鲁棒性。

通信方面，更密采样后的最小余量约为 \MarginHalf{} dB。若对既有传播损耗统一增加额外损耗，在固定保障关系与该采样集合上，能被保证承受的增量最多就是该最小值。该量很小，因此论文只报告确定性标准算例通过，不把通信覆盖外推至新的遮挡、设备偏差或天气条件。若用于实际任务，应优先重新规划最紧链路附近的中继位置和窗口，再讨论调度收益。
\FloatBarrier
